\begin{table}
\begin{tcolorbox}[tabularx={X|p{2cm}|X|},enhanced,fonttitle=\bfseries\large,fontupper=\normalsize\sffamily,
colback=yellow!10!white,colframe=red!50!black,colbacktitle=blue!30!white,
coltitle=black,title=Graph Theoretic Activities (Part 2)]
\hline
Command & Arguments & Purpose \\
\hline
\hline
\verb'-automorphism_group'  &     & Compute the automorphism group and write a report. See Section~\ref{sec:canonical:form:objects:graphs}. \\
\hline
\verb'-automorphism_group_colored_graph'  &     & Compute the automorphism group of a vertex colored graph and write a report. See Section~\ref{sec:canonical:form:objects:graphs}. \\
\hline
\verb'-properties'  &     & Compute properties of the graph. \\
\hline
\verb'-eigenvalues'  &     & Compute the eigenvalues of the graph (ordinary and Laplace) and write to a csv file. \\
\hline
\verb'-eigenvalues_report'  &     & Compute the eigenvalues of the graph (ordinary and Laplace) and produce a latex report. \\
\hline
\verb'-draw'  &  options   & Draw the graph. Uses a draw options object as in Table~\ref{tab:draw:options} for drawing. \\
\hline
\verb'-create_distance_poset'  &  v   & Create the distance poset with respect to the vertex $v$. \\
\hline
\verb'-test_if_distance_regular'  &  v   & Test if distance regular with respect to the vertex $v$. This is meaningful only if the graph is vertex transitive. \\
\hline
\end{tcolorbox}
\caption{\label{tab:graphactivities:2}Graph Theoretic Activities (Part 2)}
\end{table}
